\documentclass[11pt]{article}
\usepackage[a4paper, margin=1in]{geometry}
\usepackage[utf8]{inputenc}
\usepackage{fancyhdr}
\usepackage{hyperref}
\usepackage[hidelinks]{hyperref}

\pagestyle{fancy}
\fancyhf{}
\rhead{\thepage}
\lhead{Cover Letter}

\hypersetup{
    pdftitle={Cover Letter - DQI for AOR},
    pdfauthor={Francisco Javier Cisneros-Ruiz}
}

\begin{document}

\noindent
\textbf{Francisco Javier Cisneros-Ruiz} \\
Postgraduate Program, Instituto Mexicano del Petroleo \\
Eje Central Lazaro Cardenas Norte 152 \\
Mexico City, Mexico \\
\href{mailto:fcisnerosr@outlook.es}{fcisnerosr@outlook.es} \\
\vspace{0.5cm}

\today

\vspace{0.5cm}

\noindent
\textbf{To the Editor-in-Chief} \\
\textit{Applied Ocean Research}

\vspace{0.5cm}

Dear Editor,

\vspace{0.3cm}

We submit the enclosed manuscript titled ``Detection Quality Index (DQI): Inverse Optimisation Methodology Based on Genetic Algorithms for Damage Identification in Jacket-Type Offshore Platforms'' for consideration for publication in \textit{Applied Ocean Research}.

This work addresses a critical gap in offshore structural health monitoring: the lack of unified objective functions that simultaneously fuse multiple vibration-based indicators while penalising false positives—a key challenge in noisy operational environments. Our manuscript presents three major contributions:

\begin{enumerate}
    \item \textbf{Detection Quality Index (DQI):} A novel composite metric that combines eight modal and flexibility-based damage indicators via genetic algorithms, exponentially penalising false alarms and achieving detection rates of 92.2\% (denting) and 96.5\% (corrosion) at 30\% damage severity on a 4-leg jacket platform with 120 structural elements.

    \item \textbf{Topology-Weighted Precision:} A methodological innovation that corrects systematic underestimation of localiser performance in lattice structures by weighting false positives according to their topological distance in the structural connectivity graph, yielding up to +22 percentage point improvements for brace elements.

    \item \textbf{Exhaustive Numerical Validation:} Statistical validation over 3{,}240 simulation runs with noise robustness analysis (SNR 5--50\,dB) demonstrating applicability to real MEMS sensor configurations, including sensitivity analysis confirming invariance of structural conclusions to parametric choices.
\end{enumerate}

The methodology is directly applicable to the aging Gulf of Mexico infrastructure and represents a tangible advance toward automated condition monitoring in offshore operations. The work extends the GA-based DI combination framework of Ervin \& Zeng (2022) to the offshore jacket domain whilst introducing novel penalisation and localisation metrics adapted to the specific challenges of redundant lattice structures.

We confirm that this manuscript has not been published previously and is not under consideration for publication elsewhere. All authors concur with the submission and have no competing interests.

We would be grateful if you would consider this manuscript for publication in \textit{Applied Ocean Research}.

\vspace{0.5cm}

Sincerely,

\vspace{1.5cm}

\noindent
\textbf{Francisco Javier Cisneros-Ruiz} \\
On behalf of all authors

\vspace{1cm}

\hrule
\vspace{0.3cm}

\noindent
\textbf{Corresponding Author:} Francisco Javier Cisneros-Ruiz \\
\textbf{Email:} \href{mailto:fcisnerosr@outlook.es}{fcisnerosr@outlook.es} \\
\textbf{Institution:} Postgraduate Program, Instituto Mexicano del Petroleo \\
\textbf{Address:} Eje Central Lazaro Cardenas Norte 152, Mexico City, Mexico

\end{document}
